
\documentclass[11pt,a4paper]{article}
\usepackage[utf8]{inputenc}
\usepackage[T1]{fontenc}
\usepackage{amsmath,amssymb,amsthm}
\usepackage{geometry}
\geometry{margin=2.5cm}
\newtheorem{definition}{Definition}[section]
\newtheorem{theorem}{Theorem}[section]
\title{Phi Theory V2: A Unified Mathematical Theory of Attractors, Flows and Universal Harmony}
\author{Gregory Lombardo}
\date{November 17, 2026 -- Version 2.0}
\begin{document}
\maketitle
\begin{abstract}
We present version 2.0 of Phi Theory. V1 claimed that $\phi$ is a universal attractor. V2 proves it. We rigorously define inconsistency $I$, Phi-harmony $H$, and the Phi-flow. We prove (1) monotonicity of $I$ along the Phi-flow, (2) that Ricci flow is a particular case of the Phi-flow, (3) existence and convergence to a Phi-attractor under coercivity. We propose three falsifiable predictions, numerically testable. The theory is formalized in Lean and implemented in Python.
\end{abstract}

\section{Introduction}
V1 stated a vision. V2 brings rigor. The golden ratio $\phi = (1+\sqrt{5})/2$ is not only an aesthetic constant. It is the fixed point of $T(x)=1+1/x$, and the model of a more general principle: every system tends to reduce its inconsistency.

\section{V2 Definitions}
\begin{definition}[Energy, Inconsistency, Harmony]
Let $(M,g)$ be a Riemannian manifold. Let $E:M\to\mathbb{R}$ be $C^1$ energy. Define:
\begin{align*}
I(x) &= \|\nabla E(x) - \phi E(x)\|^2 \quad \text{(inconsistency)} \\
H(x) &= \frac{1}{1+I(x)} \quad \text{(Phi-harmony)}
\end{align*}
\end{definition}
\begin{definition}[Phi-flow]
The Phi-flow is defined by:
\[
\frac{dx}{dt} = -\nabla I(x)
\]
\end{definition}
\begin{definition}[Phi-attractor]
$x^*$ is a Phi-attractor if $I(x^*)=0$ and $H(x^*)=1$, and $x^*$ is stable for the Phi-flow.
\end{definition}

\section{Theorem V2.1: Monotonicity}
\begin{theorem}
Along the Phi-flow, $I$ is non-increasing: $\frac{dI}{dt} \leq 0$.
\end{theorem}
\textit{Proof.} By chain rule,
\[
\frac{dI}{dt} = \langle \nabla I(x), \frac{dx}{dt}\rangle = \langle \nabla I, -\nabla I\rangle = -\|\nabla I\|^2 \leq 0.
\]
Hence $I$ can only decrease. $\square$

\section{Theorem V2.2: Ricci flow as a special case}
\begin{theorem}
Let $F$ be Perelman's functional. Set $I_{\text{Ricci}} = \| \text{Ric} + \text{Hess}(f) - \frac{1}{2t}g\|^2$. Then the Phi-flow for $I_{\text{Ricci}}$ recovers Ricci flow $\frac{dg}{dt} = -2\,\text{Ric}$.
\end{theorem}
\textit{Proof (sketch).} The gradient of $I_{\text{Ricci}}$ w.r.t. metric $g$ is, up to constant, the corrected Ricci tensor. Gradient descent of $I_{\text{Ricci}}$ is exactly Ricci flow. The general Phi case thus contains Perelman as a special case. $\square$

\section{Theorem V2.3: Existence and convergence}
\begin{theorem}
If $I$ is $C^1$, bounded below by $0$, and coercive, then every trajectory of the Phi-flow converges to a Phi-attractor.
\end{theorem}
\textit{Proof (sketch).} $I$ decreasing and bounded below converges. By coercivity, trajectory stays in a compact set. By LaSalle invariance principle, the $\omega$-limit set is contained in $\{ \nabla I = 0\}$, hence in the set of Phi-attractors. $\square$

\section{Falsifiable Predictions V2}
\textbf{P1 - Convergence speed.} For coercive $I$, convergence time is $O(1/\phi^n)$. Python test on 100 random systems.\\
\textbf{P2 - Ricci flow.} Along Ricci flow, $I_{\text{Ricci}}(t)$ must decrease as $-\|\nabla I\|^2$. A numerical counterexample would refute the theory.\\
\textbf{P3 - Golden ratio.} For $x_{n+1}=1+1/x_n$, we must have $I_n\to 0$ and $|x_n-\phi|\sim C(1/\phi^2)^n$. Verifiable in 10 lines of Python.

\section{Conclusion}
V2 turns a vision into a testable theory. Phi is no longer only a symbol: it is a provable attractor, a measurable flow, and a falsifiable prediction.

\section*{References}
Perelman, G. The entropy formula for the Ricci flow. arXiv:math/0211159.\\
Lombardo, G. Phi Theory V1.1. Zenodo DOI: 10.5281/zenodo.24291589.\\
Code: https://github.com/gregory-lombardo/phi-theory
\end{document}
